% The gradient is perpendicular to the level set and points uphill.
%
%   figures/tikz/ra-gradient.tex
%       -> media/figures/ra-gradient.{png,svg}
%
% Lecture 13 (Real Analysis Review), Derivatives.
%
% GEOMETRY (1 unit = 1.3 cm). phi(x) = (x_1^2 + 4 x_2^2)/2, grad = (x_1, 4 x_2).
% Level curves x_1^2 + 4 x_2^2 = 2c are ellipses with semi-axes sqrt(2c) and
% sqrt(2c)/2, drawn for phi = 0.3, 0.8, 1.22, 1.8, 2.5. The point x = (1.2,0.5)
% has phi = (1.44 + 1.0)/2 = 1.22 and grad = (1.2,2.0). The tangent direction
% (2.0,-1.2) satisfies (1.2)(2.0) + (2.0)(-1.2) = 0, perpendicular to the
% gradient by construction. The gradient is drawn at 0.6 of its length.
%
% GREYSCALE: contours are thin solid with their values printed; the gradient is a
% solid arrow, the descent direction a dashed arrow, the tangent line dotted.
\definecolor{okabeblue}{HTML}{0072B2}

\begin{tikzpicture}[
    x=1.3cm, y=1.3cm,
    >={Latex[length=2.0mm]},
    font=\small,
    lab/.style={font=\scriptsize, inner sep=1.5pt},
  ]
  \foreach \c in {0.3,0.8,1.22,1.8,2.5} {
    \pgfmathsetmacro{\a}{sqrt(2*\c)}
    \pgfmathsetmacro{\b}{sqrt(2*\c)/2}
    \draw[black!60, line width=0.5pt] (0,0) ellipse (\a cm*1.3 and \b cm*1.3);
  }
  \draw[black!35, ->] (-2.3,0) -- (2.5,0) node[right, lab] {$x_1$};
  \draw[black!35, ->] (0,-1.4) -- (0,1.5) node[above, lab] {$x_2$};
  % highlight the level curve through x
  \draw[okabeblue, line width=1.1pt] (0,0) ellipse (1.5620 cm*1.3 and 0.7810 cm*1.3);
  \filldraw (1.2,0.5) circle (1.8pt);
  % tangent line, direction (2.0,-1.2)/2.332
  \draw[densely dotted, thick] (1.2,0.5) ++(-0.6,0.36) -- ++(1.2,-0.72);
  % gradient (1.2,2.0)*0.4 and negative gradient
  \draw[->, very thick] (1.2,0.5) -- ++(0.48,0.8);
  \draw[->, thick, dashed] (1.2,0.5) -- ++(-0.48,-0.8);
  % right-angle mark between gradient direction (0.514,0.857) and tangent (0.857,-0.514)
  \draw[black!70, line width=0.4pt] (1.2,0.5) ++(0.1028,0.1714) -- ++(0.1714,-0.1028) -- ++(-0.1028,-0.1714);
  \node[lab, anchor=south east] at (1.62,1.28) {$\nabla\phi(\mathbf{x})$};
  \node[lab, anchor=north east, fill=white, inner sep=1pt] at (0.8,-0.25) {$-\nabla\phi(\mathbf{x})$};
  \node[lab, anchor=north west, fill=white, inner sep=0.5pt] at (1.27,0.42) {$\mathbf{x}$};
  \node[lab, anchor=west, okabeblue, fill=white, inner sep=1pt] at (1.62,-0.6) {$\phi = 1.22$};
  \node[lab, anchor=west, fill=white, inner sep=1pt] at (2.04,-0.2) {$\phi = 1.8$};
\end{tikzpicture}
